\section{Proof of distributional alignment}
\label{sec:alignment-proof}

We prove Lemma~\ref{lem:alignment}, the combinatorial input used to
choose the sampling law on layer sets. The argument depends only on
the list bound and the order of the indices; it has no graph-theoretic
or analytic hypotheses.

\begin{proof}[Proof of Lemma~\ref{lem:alignment}]
Call $B$ \emph{bad} if no map in \eqref{align:conclusion} exists.  For each
label $x$ that appears on a list indexed inside $B$, consider
\[
 \bigcap\{I:\varnothing\ne I\subseteq B,\ x\in A_I\}.
\]
The set $B$ is bad exactly when one of these intersections is empty.
Indeed, a nonempty intersection permits a choice of $a(x)$ independently
for each appearing label; labels absent from all these lists may be sent to
any element of the nonempty set $B$.

We first prove that for every $\xi>0$, some $r\geq s$ admits a real
probability distribution on $\binom{[r]}s$ under which every list family
satisfying \eqref{align:separation} has bad-set probability at most $\xi$.
The proof has three steps.  Finite minimax turns a failure of this assertion
into random list patterns making every $s$-set bad with positive
probability.  Ramsey homogenization produces an order-invariant limiting
pattern on the rationals.  Finally, an invariant-cut argument shows that
such a limiting pattern has no bad sets.

\paragraph{Finite patterns and minimax.}
Order each list arbitrarily.  On a finite ordered index set, retain only
the length of each list and the equality relation among its occupied
slots.  Each list has at most $d$ slots, and no two occupied slots of one
list are equal.  There are finitely many resulting patterns.  Every actual
list family gives such a pattern, and conversely a pattern is realized by
using its equivalence classes as labels.  Thus the original label universe
has no further role in either separation or badness.

Write $\mathcal C_r$ for the finite set of patterns on $[r]$ satisfying
\eqref{align:separation}, and put
\[
 b(B,C)=\boldsymbol{1}\{B\text{ is bad in }C\},
 \qquad B\in\binom{[r]}s,\ C\in\mathcal C_r.
\]
If the assertion with tolerance $\xi$ were false for every $r$, then for
each $r\geq s$ the finite minimax theorem \cite{vonNeumann1928} would give
\[
 \min_{\nu}\max_{C\in\mathcal C_r}
       \E_{B\sim\nu}b(B,C)
 =
 \max_{\rho}\min_{B\in\binom{[r]}s}
       \E_{C\sim\rho}b(B,C)
 >\xi.
\]
Here $\nu$ and $\rho$ range over the two finite probability simplices;
their compactness ensures that the extrema are attained.  Consequently,
for every $r\geq s$ there would be a random allowed pattern $C_r$ such that
\begin{equation}\label{align:uniform-bad}
 \Pr(B\text{ is bad in }C_r)\geq\xi
 \qquad\text{for every }B\in\binom{[r]}s.
\end{equation}
Assume this counterstatement for the remainder of the argument.

\paragraph{An order-invariant limit.}
The homogenization step applies finite Ramsey theory \cite{Ramsey1930}
to discretized probability patterns, following the kind of extraction
used by Trotter and Winkler \cite{TrotterWinkler1998}. We give the
particular argument needed here in full.
For $t\geq1$, let $\mathcal P_t$ denote the finite set of all allowed
patterns on $[t]$, using lists on subsets of size at most $s$.
Restriction to an ordered subset, followed by its increasing identification
with an initial interval of integers, defines a map between the appropriate
pattern spaces.

Fix $k\geq s$.  For each $1\leq t\leq k$ and each $t$-subset $J$ of
$[r]$, the restriction of $C_r$ to $J$ has a probability vector in the
finite simplex on $\mathcal P_t$.  Color $J$ by this vector with every
entry placed in a bin of width at most $1/k$.  The number of colors depends
only on $s,d,t,k$, not on $r$.  Successive applications of the finite
Ramsey theorem, with the required intermediate set sizes chosen backwards,
give the following when $r$ is sufficiently large: there is a $k$-element
set $H_k$ on which all the colorings, for $1\leq t\leq k$, are
homogeneous.  Passing to a subset preserves each homogeneity already
obtained.  In particular, the induced pattern laws on any two $t$-subsets
of $H_k$ differ by at most $1/k$ in every coordinate.

Let $P_t^{(k)}$ be the law on the first $t$ points of $H_k$.  If
$t\leq u\leq k$ and $J\in\binom{[u]}t$, restricting $P_u^{(k)}$ to
$J$ gives the law on the corresponding $t$-subset of $H_k$.  Therefore
\begin{equation}\label{align:approx-consistency}
 \left|
   \bigl(P_u^{(k)}|_J\bigr)(C)-P_t^{(k)}(C)
 \right|\leq\frac1k
 \qquad(C\in\mathcal P_t).
\end{equation}
Compactness of each finite simplex and a diagonal subsequence as
$k\to\infty$ produce limiting laws $P_t$ for every $t$.  Taking limits
in \eqref{align:approx-consistency} gives exact consistency under every
ordered restriction:
\begin{equation}\label{align:consistency}
 P_u|_J=P_t
 \qquad(t\leq u,\ J\in\binom{[u]}t).
\end{equation}
Moreover, \eqref{align:uniform-bad} gives
\begin{equation}\label{align:limit-bad}
 P_s(\text{the full index set is bad})\geq\xi.
\end{equation}

Assign the law $P_t$ to every increasing $t$-tuple of rational indices.
The consistency in \eqref{align:consistency} constructs a countable random
array of list lengths and equality relations indexed by the nonempty
$I\subseteq\Q$ of size at most $s$.  For completeness, enumerate $\Q$
and successively extend the pattern on the first $n$ enumerated indices
to the first $n+1$, using the conditional probabilities supplied by their
consistent finite laws.  The choices on zero-probability conditioning
events are immaterial.  Every finite set of rational indices then has its
specified law.

All equivalence-relation axioms, distinctness within a list, and separated
disjointness hold simultaneously almost surely: each violation involves
finitely many indices, and there are only countably many possible
violations.  The array's law is invariant under every increasing
automorphism of $\Q$, because its finite laws depend only on order.
Equation~\eqref{align:limit-bad} says that every fixed rational $s$-set
still has bad probability at least $\xi$.

We have thus retained the positive bad-set probability while removing all
distinctions between index sets having the same order type.  We next show
that this order invariance, together with separated disjointness, forces
every label to have an index common to all its occurrences.

\paragraph{An invariant cut for each label.}
Fix a nonempty $I\subseteq\Q$, $|I|\leq s$, and one of its $d$ possible
slots.  Let $E$ be the event that this slot is occupied.  On $E$, write
$\mathcal L$ for its equivalence class of occupied slots, and define
\begin{equation}\label{align:cut}
 p=\sup\{\min J:\text{a slot at }J\text{ belongs to }\mathcal L\}.
\end{equation}
The set whose supremum is taken contains $\min I$.  It is bounded above
by $\max I$, since an occurrence at $J$ with $\min J>\max I$ would
violate separated disjointness.  Thus
\begin{equation}\label{align:cut-hull}
 \min I\leq p\leq\max I
 \qquad\text{on }E.
\end{equation}
The random variable $p$ is measurable on $E$: the occurrence family is
countable, and testing whether its supremum is at most a given real number
is a countable intersection of measurable conditions on slots.

Every increasing automorphism $g$ of $\Q$ extends uniquely to an increasing
homeomorphism $\bar g$ of $\R$.  Transporting the array by $g$ transports
the occurrence family in \eqref{align:cut} and sends its supremum to
$\bar g(p)$.  If $g$ fixes $I$ pointwise, it also fixes the selected slot
and its occupancy event.  Hence the finite measure
\[
 \nu(U)=\Pr(E\text{ and }p\in U),
 \qquad U\subseteq\R\text{ Borel},
\]
is invariant under all such $\bar g$.  This formulation also covers
$\Pr(E)=0$ without conditioning on a null event.

Let $a<b$ be consecutive points of $I$.  For any rational
$a<q<q'<b$, there is an increasing automorphism of $\Q$ fixing $I$
pointwise and carrying $q$ to $q'$.  For example, use a piecewise linear
increasing bijection of $\R$ that is the identity off $[a,b]$, maps $q$
to $q'$, and is linear on $[a,q]$ and $[q,b]$.  Its rational breakpoints
and rational slopes make it a bijection of $\Q$.  Invariance gives
\[
 \nu(({-\infty},q])=\nu(({-\infty},q']),
 \qquad\text{so}\qquad \nu((q,q'])=0.
\]
Countably many such rational intervals cover $(a,b)$, whence
$\nu((a,b))=0$.  Together with \eqref{align:cut-hull}, this proves
\begin{equation}\label{align:cut-index}
 \Pr(E\text{ and }p\notin I)=0.
\end{equation}
If $I$ is a singleton, \eqref{align:cut-index} follows directly from
\eqref{align:cut-hull}.

There are only countably many possible slots, so
\eqref{align:cut-index} holds simultaneously at every occupied slot with
probability one.  Equivalent slots have exactly the same occurrence
family and hence the same supremum $p$.  It follows that every equivalence
class has a common index in all the sets on whose lists it occurs.
In particular, on every finite rational $s$-set $B$, every label appearing
inside $B$ has a nonempty intersection of its occurrence sets there.
Thus no such $B$ is bad, contradicting \eqref{align:limit-bad} and
$\xi>0$.  This proves the assertion with real probabilities.

\paragraph{Rational probabilities.}
Choose $0<\xi<\eta$ and obtain $r$ and a real distribution $\mu_0$ with
bad-set probability at most $\xi$ for every allowed pattern.  Approximate
$\mu_0$ by a rational point $\mu$ of the same finite simplex with
\[
 \sum_{B\in\binom{[r]}s}|\mu(B)-\mu_0(B)|<\eta-\xi.
\]
Every bad-set indicator takes values in $[0,1]$, so this changes its
expectation by less than $\eta-\xi$, uniformly over all patterns.
The distribution $\mu$ therefore has the required guarantee.
\end{proof}

\begin{remark}\label{align:effective}
For rational $\eta>0$, the finite data in Lemma~\ref{lem:alignment} can
also be found by a terminating search.  Enumerate $r\geq s$, enumerate
its finitely many allowed slot patterns, and minimize their maximum
bad-set probability by a rational linear program.  Search until its
optimum is less than $\eta$.  The proof with a smaller tolerance guarantees
termination.  No bound on the label universe, or oracle describing it, is
needed.
\end{remark}
